% GPT-6 Astra Pro; effort not exposed; /root/replacement; 22 November 2026.
\section[HD-F1-R33: two forward maps are impossible]{HD-F1-R33: the common-quotient branch is excluded}
\label{sec:followup-rm-two-forward-maps-exclusion}

There is no actual object in the rank-$(3,2)$,
equal-slope semistable TF adjacent gate of
R21--R28 with
\[
                 \dim\operatorname{Hom}_G(A,T)=2.
 \tag{F1R33.1}
\]
The genuine real-quadratic fourfold, the specified
coherent faithful finite point action and actual
lift, gluability, weak semiregularity and the
whole balanced character are exactly those of
that gate. In particular this excludes the
entire common-quotient branch of R22--R32,
including its remaining actual surface
extensions. It also excludes the $h=2$
common-image branch. At these endpoint ranks
the only branch left by this result is the
common-image case $h=3$; no object in that
case is asserted.

The contradiction is
\[
 \operatorname{ext}^2_G(E,E)
 \ \ge\ \operatorname{ext}^2_G(A,A)
          +\operatorname{ext}^2_G(T,T)-4
 \ \ge\ 20\ >\ 18.
 \tag{F1R33.2}
\]
We prove both inequalities through the
original character and triangle maps.

\paragraph{1. The actual triangle and its whole characters.}
Perform only the actual determinant line
twist already specified in R21. It gives
the actual lifted triangle
\[
 A\xrightarrow{i}E\xrightarrow{p}C=T[-1]
       \xrightarrow{d}A[1],\qquad d=\delta[-1],
       \qquad \delta:T\longrightarrow A[2],
 \tag{F1R33.3}
\]
with $\operatorname{ch}(E)=v_q=v_1v_2$,
$v_i=1-qA_i$, $q>0$. Here and throughout,
$\operatorname{Ext}_G$ means the invariant
ordinary Ext space of the specified coherent
lifts. The finite-action averaging and
identity-sector comparisons are the ones
already proved for these exact lifts.

Put $W=\operatorname{Hom}_G(A,T)$ and assume
for the contradiction that $\dim W=2$.
The full character results R26--R28 give
\[
 \begin{split}
 z_T:=\operatorname{ch}(T)&=a_0v_1v_2+c\,p_1p_2,\\
 z_A:=\operatorname{ch}(A)&=(a_0+1)v_1v_2+c\,p_1p_2,
 \end{split}
 \qquad
 a_0>0,\quad c\ne0,\quad p_i\notin\mathbf C v_i.
 \tag{F1R33.4}
\]
These are the whole characters of the actual
endpoints, with R26's rationality and
factor-exchanging field maps retained.
The inequalities $a_0>0$ and $c\ne0$
are proved in R27--R28 from their actual
effective target cycles. No factor $p_i$
is being realized by a proposed sheaf.

Translations act trivially on the invariant
Hochschild representatives, and a line tensor
from $\operatorname{Pic}^0(X)$ has zero first
Chern class in the characteristic conjugation
formula. Thus the specified point action
fixes $HH^2(X)$. Naturality of the actual
characteristic transformation at each chosen
lift makes its values invariant. In particular,
for $M=A,T$ the full maps factor as
\[
 c_{z_M}^2:
 HH^2(X)\xrightarrow{\operatorname{ev}_M^2}
       \operatorname{Ext}^2_G(M,M)
       \xrightarrow{\sigma_M^2}HH_{-2}(X).
 \tag{F1R33.5}
\]
Consequently the rank of this full
character map is a lower bound for the
indicated invariant Ext dimension.

\paragraph{2. The two full character ranks sum to at least twenty-four.}
Use the actual cyclic exterior isomorphism
of R19 and R26. Write
\[
 U=U_1\oplus U_2,\qquad \dim U_i=4,\qquad
 \Phi_i:\bigwedge^2U_i\xrightarrow{\sim}S_i^0,
 \quad \xi_i=\Phi_i^{-1}(v_i),\quad
 \alpha_i=\Phi_i^{-1}(p_i).
 \tag{F1R33.6}
\]
Under the full module isomorphism, the
degree-two character map is multiplication
by the corresponding four-vector from
$\bigwedge^2U$ to $\bigwedge^6U$.
The two four-vectors here are
\[
 w_T=a_0\xi_1\xi_2+c\alpha_1\alpha_2,
 \qquad
 w_A=(a_0+1)\xi_1\xi_2+c\alpha_1\alpha_2.
 \tag{F1R33.7}
\]
Products are exterior products. Each pair
$\xi_i,\alpha_i$ is linearly independent
by (F1R33.4).

Both multiplication maps preserve the
three indicated source and target blocks:
\[
 \begin{array}{c|c}
 \text{source}&\text{target}\\ \hline
 \bigwedge^2U_1&\bigwedge^4U_1\otimes\bigwedge^2U_2\\
 U_1\otimes U_2&\bigwedge^3U_1\otimes\bigwedge^3U_2\\
 \bigwedge^2U_2&\bigwedge^2U_1\otimes\bigwedge^4U_2.
 \end{array}
 \tag{F1R33.8}
\]
Fix volume forms in the two blocks.
For a nonzero coefficient $t$ equal to
$a_0$ or $a_0+1$, the first block map is
\[
 k\longmapsto t[k\xi_1]\xi_2+c[k\alpha_1]\alpha_2.
 \tag{F1R33.9}
\]
The exterior pairing on $\bigwedge^2U_1$
is nondegenerate. Thus the two coefficient
functionals in (F1R33.9) are independent.
The two target vectors are independent
as well, and both coefficients $t,c$ are
nonzero. This block has rank two.
The third block has rank two by the same
argument. Their targets are distinct,
so each full character map already has
rank four on these two blocks.

Let $B_T,B_A$ be the two maps on the
middle, sixteen-dimensional source block.
Their difference is exactly
\[
 B_A-B_T:
       u_1\otimes u_2\longmapsto
                  (u_1\xi_1)\otimes(u_2\xi_2).
 \tag{F1R33.10}
\]
This is an isomorphism. Indeed, in R26's
specified basis,
\[
       \xi_i=-t_{i1}t_{i2}-q w_{i1}w_{i2}.
\]
Multiplication by $\xi_i$ sends the four
basis vectors to nonzero scalar multiples
of the four distinct complementary
three-vectors; $q\ne0$. It therefore gives
an isomorphism $U_i\to\bigwedge^3U_i$.
The two degree-two factors commute across
the block reordering, so the sign in
(F1R33.10) is positive with the retained
block order.

The actual linear map
\[
 U_1\otimes U_2\longrightarrow
          \operatorname{im}B_A\oplus\operatorname{im}B_T,
                 \quad x\longmapsto(B_Ax,B_Tx)
 \tag{F1R33.11}
\]
is injective: its kernel is killed by
the isomorphism (F1R33.10).
It follows that
$\operatorname{rank}B_A+\operatorname{rank}B_T\ge16$.
Adding the four even-block ranks for each
endpoint proves
\[
       \operatorname{rank}c_{z_A}^2
           +\operatorname{rank}c_{z_T}^2\ge24,
 \qquad
       \operatorname{ext}^2_G(A,A)
           +\operatorname{ext}^2_G(T,T)\ge24.
 \tag{F1R33.12}
\]
This is a rank estimate for the full
twenty-eight-dimensional degree-two
source. It is not the rank-two estimate
for its balanced-kernel restriction.

\paragraph{3. The original coupling has a nondegenerate two-dimensional trace pairing.}
Let $K=K_1\oplus K_2$ be the balanced
degree-two kernel. Weak semiregularity
and the exact endpoint Hom sequences
give the original Yoneda memberships
\[
 \operatorname{ev}_A^2(K)\subset\nu(W),\qquad
 \operatorname{ev}_T^2(K)\subset\mu(W),\qquad
 \nu(f)=\delta f,\quad \mu(f)=f[2]\delta.
 \tag{F1R33.13}
\]
R26 gives rank two for the full traced
image of each restriction. Since $W$
has dimension two, $\sigma_A^2\nu$
and $\sigma_T^2\mu$ are injective.
In particular $\nu$ and $\mu$ are
injective. Full cyclic Atiyah trace gives
\[
                \sigma_A^2(\delta f)
                      =\sigma_T^2(f[2]\delta).
 \tag{F1R33.14}
\]
The moved degrees are zero and two,
so the interchange sign is positive.
Since $z_A-z_T=v_q$ and $K v_q=0$,
injectivity makes the two coefficients
in (F1R33.13) identical. There is an
actual surjective linear map
\[
 f_K:K\twoheadrightarrow W,\qquad
 \operatorname{ev}_A^2(k)=\delta f_K(k),\qquad
 \operatorname{ev}_T^2(k)=f_K(k)[2]\delta.
 \tag{F1R33.15}
\]
This argument applies to either
rank-one factorization of R21.

Choose $k_i\in K_i$ with
$D_{k_i}p_i=\omega_i\ne0$, rescaling
the nonzero functional of R26.
Here $\omega_i$ spans the pure block
$H^{0,2}_i$. Their characteristic
images have distinct block degrees,
so $f_i=f_K(k_i)$ form a basis of $W$.
Put $m_i=\mu(f_i)=\operatorname{ev}_T^2(k_i)$.
The Buchweitz--Flenner module identity,
followed by its actual decomposition
map and the cyclic exterior action,
gives
\[
 \begin{split}
 \operatorname{Tr}_T(m_i[2]m_i)
                         &=D_{k_i}^2z_T=0,\\
 \operatorname{Tr}_T(m_1[2]m_2)
                         &=D_{k_1}D_{k_2}z_T
                           =c\,\omega_1\omega_2\ne0.
 \end{split}
 \tag{F1R33.16}
\]
These equalities take values in
$H^4(X,\mathcal O_X)$. In degree four
this is the sole full-trace component
on the fourfold, so the full trace is
the ordinary Serre trace. The square
vanishes because a second action in
one block exceeds that block's
antiholomorphic degree two. In the
cross product, $K_i v_i=0$ removes
the balanced term and each $p_i$
contributes its nonzero $\omega_i$.
Both operators have even degree,
so their interchange again has sign
positive. The formula uses the actual
Hochschild module action, not an
assumed algebra property of an
arbitrary HKR isomorphism.

Fix the invariant volume-form
trivialization of the canonical
bundle. Ordinary Serre duality is
equivariant for the point action;
restriction to invariants remains
nondegenerate by finite averaging.
Thus
\[
           \operatorname{Ext}^4_G(T,A)\simeq W^*.
 \tag{F1R33.17}
\]
Consider the original triple-composition map
\[
 J_\delta:W\longrightarrow\operatorname{Ext}^4_G(T,A),
             \quad f\longmapsto\delta[2]f[2]\delta.
 \tag{F1R33.18}
\]
Its pairing with $g\in W$ is exactly
\[
 \operatorname{Tr}_T\bigl(g[4]\delta[2]f[2]\delta\bigr)
              =\operatorname{Tr}_T\bigl(\mu(g)[2]\mu(f)\bigr).
 \tag{F1R33.19}
\]
Equations (F1R33.16)--(F1R33.19)
give its matrix, after a fixed
nonzero volume scalar, as
\[
                         \begin{pmatrix}0&c'\\c'&0\end{pmatrix},
                                  \qquad c'\ne0.
 \tag{F1R33.20}
\]
It follows that $J_\delta$ is an
isomorphism. All maps in this pairing
use the same actual $\delta$.

\paragraph{4. Actual morphisms of the triangle give the Ext bound.}
Define the compatibility map
\[
 \begin{split}
 D:\operatorname{Ext}^2_G(A,A)\oplus
                 \operatorname{Ext}^2_G(T,T)
        &\longrightarrow\operatorname{Ext}^4_G(T,A),\\
 (a,t)&\longmapsto a[2]\delta-\delta[2]t.
 \end{split}
 \tag{F1R33.21}
\]
It is surjective: taking $(a,t)=(\nu(f),0)$
gives exactly $J_\delta(f)$. Thus
\[
 \dim\ker D=\operatorname{ext}^2_G(A,A)
                    +\operatorname{ext}^2_G(T,T)-2.
 \tag{F1R33.22}
\]

Let $\mathcal M$ be the vector space
of actual morphisms from (F1R33.3)
to its shift by two. Concretely it
consists of triples $(a,e,t)$ satisfying
\[
 \begin{gathered}
 a:A\to A[2],\quad e:E\to E[2],\quad t:T\to T[2],\\
 e i=i[2]a,\qquad p[2]e=t[-1]p,
                  \qquad a[2]\delta=\delta[2]t.
 \end{gathered}
 \tag{F1R33.23}
\]
All are morphisms for the given lifts.
Projection to $(a,t)$ maps $\mathcal M$
onto $\ker D$. To verify surjectivity,
rotate the original distinguished
triangle twice to
\[
 C\xrightarrow{d}A[1]\xrightarrow{-i[1]}E[1]
                         \xrightarrow{-p[1]}C[1].
 \tag{F1R33.24}
\]
The given pair $(t[-1],a[1])$ commutes
with its first arrow precisely when
$D(a,t)=0$. Axiom TR3 supplies a
morphism on $E[1]$. Shift it back
to obtain $e$. The two minus signs
in (F1R33.24) occur in both rows
and cancel, giving exactly the first
two equations in (F1R33.23).
The shift by two changes no triangle
arrow sign. Equivalently one can
fill in the ordinary category and
average the filler: all the displayed
conditions are linear and their
specified endpoints are invariant.
No functorial choice of a filler
is needed for this surjectivity of
finite-dimensional vector spaces.

Project instead to the middle map
$e\in\operatorname{Ext}^2_G(E,E)$.
We determine this projection's kernel
exactly. At $e=0$, the two actual
Hom sequences of (F1R33.3) give
\[
 \begin{split}
 \ker\bigl[\operatorname{Ext}^2_G(A,A)
                 \xrightarrow{i[2]\circ(-)}
                        \operatorname{Hom}_G(A,E[2])\bigr]
                         &=\nu(W),\\
 \ker\bigl[\operatorname{Ext}^2_G(C,C)
                 \xrightarrow{(-)p}
                        \operatorname{Hom}_G(E,C[2])\bigr]
                         &=\mu(W)[-1].
 \end{split}
 \tag{F1R33.25}
\]
For the first sequence, the preceding
group is $\operatorname{Hom}_G(A,C[1])=W$
and its map is composition with
$d[1]=\delta$. For the second,
the preceding group is
$\operatorname{Hom}_G(A[1],C[2])=W$;
its map is composition with $d$,
and shifting by one gives $f[2]\delta$.
Thus every pair in the middle
projection's kernel has uniquely
the form $a=\nu(f)$, $t=\mu(g)$,
because both Yoneda maps are injective.
Its final compatibility equation is
\[
                         J_\delta(f-g)=0.
 \tag{F1R33.26}
\]
The isomorphism (F1R33.18) forces
$f=g$. Conversely every
$(\nu(f),0,\mu(f))$ satisfies all
three equations: the first two
follow from consecutive compositions
in the original triangle, and the
third is the identical triple
composition. Therefore this kernel
is exactly the two-dimensional
image of
\[
                 W\longrightarrow\mathcal M,
                       \quad f\longmapsto(\nu(f),0,\mu(f)).
 \tag{F1R33.27}
\]

Taking dimensions of these two
actual projections gives
\[
 \begin{split}
 \operatorname{ext}^2_G(E,E)
   &\ge \dim\mathcal M-2\\
   &\ge \dim\ker D-2\\
   &=\operatorname{ext}^2_G(A,A)
             +\operatorname{ext}^2_G(T,T)-4.
 \end{split}
 \tag{F1R33.28}
\]
This proof has retained every
endpoint compatibility and every
filler ambiguity; it does not
discard a surface Ext predecessor.

\paragraph{5. The contradiction and exact closed branch.}
For the actual balanced object $E$,
weak semiregularity, invariant CY4
duality and the whole characteristic
rank eighteen give
\[
                       \operatorname{ext}^2_G(E,E)=18,
 \tag{F1R33.29}
\]
as already proved in R07 and R20.
The exact maps are the identity-sector
comparison to the ordinary full trace,
its Serre-transpose characteristic
evaluation, and their composite
$c^2_{v_q}$ of rank eighteen.
Injectivity of that full trace makes
its transpose surjective; hence
the intermediate invariant Ext space
has exactly that composite rank.
Equations (F1R33.12) and
(F1R33.28) instead give dimension
at least twenty. This proves the
asserted nonexistence.

R22 makes every actual common-quotient
object have $h=2$, so all of that
branch is now excluded, with its
original coupling and all its
coherent defects retained.
R23 bounds the common-image branch
by $h=2$ or $3$. Its first case
is likewise excluded. The proof
does not decide the common-image
case $h=3$, higher endpoint ranks,
or unrestricted amplitudes and
semistability conditions. The
previous ordinary-simplicity result
R32 is consistent with this stronger
exclusion but is not a premise of
the dimension contradiction.

\paragraph{Published framework and checking scope.}
The exterior maps are the complete
R19/R26 module isomorphisms, with
every new block and rank comparison
given in the displays (F1R33.6) through (F1R33.12).
The trace products use the actual
Hochschild module identity of
Buchweitz--Flenner, \emph{Advances
in Mathematics} 217 (2008),
243--281, DOI 10.1016/j.aim.2007.06.013,
Propositions 6.2.1 and 6.3.1 and
Corollary 6.3.2. Their complete
proofs and the graded-centre sign
were freshly read in the primary
author version, pp.~29--31.
The exact triangle definitions,
rotation signs and TR3 filler are
Stacks Project, Section 13.3,
Tag 05QK, Definitions 13.3.1--13.3.2.
The full proof that the two Hom
functors give the exact sequences
used here is Lemma 13.4.2, Tag 0149;
it was freshly read. The actual
rotated diagram and its exact
kernel are proved above.
All character operations arise
from the original triangle's
equivariant $K_0$ relation and the
previously specified determinant
line tensor. No numerical factor
is substituted for an object,
and no unrefereed theorem is a
premise of this exclusion.
